% Standalone mathematical fragment, not an end-to-end SAT theorem.
\section{A quadratic majorant for rectangular multiplication}

The following result concerns the rectangular exponent envelope supplied by
$\omega(1,k,1)\le 1+k+1/(4k)$ for $1/2\le k\le1$, together with
$\omega(1,1/2,1)\le2$. These shape bounds come from~\cite{uznanski}. The finite-algorithm
and exact-ring transfer are treated in Section~\ref{sec:matrix}; the
finite shape library needed here is specified below. The first theorem
is an exact algebraic inequality for the displayed function $U$.

For three nonnegative logarithmic dimensions, sort them as $0\le a\le b\le c$
and define
\[
 U(a,b,c)=\begin{cases}
 b+c,&2a\le b,\\
 c+a+b^2/(4a),&b\le2a,\ a>0.
 \end{cases}
\]
For $a=b=0$ set $U=c$. This is the equal-outer-dimension rectangular
curve followed by blocking the largest dimension. It is not asserted to be
the strongest general three-dimensional matrix exponent.

\begin{theorem}[Symmetric quadratic majorant]
For all $a,b,c\ge0$, with $U$ interpreted after sorting,
\[
 U(a,b,c)^2\le \frac{55}{96}
 \left(7(a^2+b^2+c^2)-4(ab+ac+bc)\right).
\tag{Q}
\]
Equality occurs at any positive scalar multiple or permutation of
$(1,1,20/17)$. In particular, $55/96$ cannot be reduced while keeping
the displayed quadratic form fixed.
\end{theorem}
\begin{proof}
If $b=0$ the result is immediate. Otherwise divide by $b^2$ and write
$x=a/b$, $y=c/b$. First suppose $1/2\le x\le1$, $y\ge1$. The difference
between the right side and $(y+x+1/(4x))^2$ is $Ay^2+By+C$, where
\[
 A=\frac{289}{96},\qquad
 B=-\frac{103x^2+55x+12}{24x},\qquad
 C=\frac{289x^2-220x+337-6/x^2}{96}.
\]
A direct rational calculation gives
\[
 4AC-B^2=\frac{(1-x)R(x)}{2304x^2},\qquad
 R(x)=-41085x^3+67815x^2-7590x-2310.
\]
On $[1/2,1]$ the degree-three Bernstein coefficients of $R$, after
$x=(1+t)/2$, are
\[
 \frac{45705}{8},\quad10615,\quad\frac{32065}{2},\quad16830.
\]
They are strictly positive. Hence $A>0$ and $4AC-B^2\ge0$, which imply
$Ay^2+By+C\ge0$ for every real $y$.

For $0\le x\le1/2$ the squared left side is $(1+y)^2$. The derivative
with respect to $x$ of the difference is
\[
 \frac{55}{96}(14x-4-4y)<0\qquad(y\ge1).
\]
The difference is therefore at least its value at $x=1/2$, already
covered. At $x=1$, the discriminant vanishes and the quadratic has its
zero at $y=20/17$, proving the equality claim.
\end{proof}

\paragraph{Uniform interpretation with positive slack.}
Fix $\delta>0$. Choose a finite rational mesh in $[1/2,1]$, containing
both endpoints, and a positive slack for each of the finitely many
algorithms from~\cite{uznanski}. For $1/2\le a/b\le1$, round $a/b$
up to the next mesh point $k$, pad the inner dimension, and block the
largest dimension. The exponent is at most
\[
 c+b\left(k+\frac1{4k}+\eta\right).
\]
The derivative of $k+1/(4k)$ lies in $[0,3/4]$ on this interval.
Thus a sufficiently fine fixed mesh and sufficiently small fixed
$\eta>0$ bound the padding and algorithmic loss by
$\delta(a+b+c)$. For $a/b<1/2$, use the half-width endpoint. Rounding
finite-base recursion counts contributes only fixed dimension factors.
The mesh can be selected using exact integer-power comparisons of the
input dimensions, so no real-number oracle is required. Proposition~\ref{prop:ring}
applies to this finite collection after taking one common number field
and the union of its finite exceptional-prime sets.

Consequently every fixed positive enlargement of a coefficient derived
from $U$ is realized by a fixed finite library. The exact envelope and its
limiting coefficient do not assert attainment of an exponent infimum
with zero slack. The main time and space coefficients use the separate
fixed-slack endpoint library of Lemma~\ref{lem:rectangle}.

\section{A final-triangle certificate from the affine plane of order three}

Identify the nine group labels with $\mathbb F_3^2$. Its four directions
partition the points into three parallel lines. For direction $j$, let
$H_j$ be the three-by-36 zero-one matrix whose rows indicate the three
unordered pairs of line classes. For a nonnegative unordered-edge entropy
vector $g\in\mathbb R_{\ge0}^{36}$, its three aggregate dimensions are
$H_jg$.

\begin{theorem}[Final triangle only]
There is a direction $j$ with
\[
 U(H_jg)\le\sqrt{\frac{1485}{128}}\,\|g\|_2.
\tag{A}
\]
In the inherited nine-group mode-entropy normalization, this gives the
local final-triangle coefficient
\[
 \beta_{9,\mathrm{final}}^2=\frac{55}{192},\qquad
 \beta_{9,\mathrm{final}}=0.5352180241\ldots.
\]
This statement does not pay for producing the three line-block tensors.
\end{theorem}
\begin{proof}
Stack the four $H_j$ into a twelve-by-36 matrix $H$. Each row has nine
ones. Within a direction the rows have disjoint supports; between two
directions any two rows have inner product two. The latter count is
obtained from the four intersections of two selected lines in each
direction: exactly the two opposite corner pairs belong to both supports.
Thus $HH^T$ has diagonal blocks $9I_3$ and off-diagonal blocks $2J_3$.
Let $K$ be block diagonal with four copies of $9I_3-2J_3$.
By (Q),
\[
 \frac14\sum_{j=1}^4 U(H_jg)^2
 \le g^T\left(\frac{55}{384}H^TKH\right)g.
\]
The nonzero eigenvalues of $H^TKH$ are $81$ with multiplicity nine and
$9$ with multiplicity three. To see this directly, decompose the
12-dimensional row space into the eight-dimensional subspace having
coordinate sum zero in each block, the three-dimensional subspace of
blockwise constants whose total is zero, and the one-dimensional constant
subspace. On these, $HH^T$ acts by $9,3,27$, respectively, while $K$
acts by $9,3,3$. Their product therefore acts by $81,9,81$.
The largest eigenvalue of the averaged quadratic form is
$55\cdot81/384=1485/128$. The minimum of four nonnegative numbers is at
most their root-mean-square, proving (A).

If the directed mode entropy vector has Euclidean norm at most
$(1+o(1))\mathcal K/9$, then joining the two orientations gives
$\|g\|_2\le(\sqrt2/9+o(1))\mathcal K$. Substitution into (A) yields
$\beta^2=2(1485/128)/81=55/192$.
\end{proof}

\paragraph{Scope of the result.}
The proof covers the final aggregation of three precomputed blocks. To
interpret it algorithmically with a fixed finite library, replace
$\beta_{9,\mathrm{final}}$ by any strictly larger coefficient, as explained
above. It gives no bound for producing those block tensors, for a complete
contraction plan, or for the full 3-SAT recurrence. The internal line-block
merges and recursive cache construction must still be charged. No verified
combination that improves the main fully paid coefficient is established
here, and this appendix's coefficient is not used in that calculation.
A common quadratic certificate for both internal and final operations,
or an adaptive combination of linear and quadratic certificates, remains
a possible continuation of this local result.
